\documentclass[12pt]{article}

\usepackage[a4paper, margin=1in]{geometry}
\usepackage{setspace}
\usepackage{booktabs}
\usepackage{graphicx}
\usepackage{hyperref}

\onehalfspacing

\title{\textbf{Model Evaluation Report}\\
Customer Churn Prediction using Random Forest Classifier}
\author{Capstone Project -- Milestone 1}
\date{}

\begin{document}

\maketitle

\section{Problem Statement}

Customer churn is a significant challenge in the telecommunications industry, as retaining existing customers is considerably more cost-effective than acquiring new ones. The objective of this project is to develop a machine learning model capable of predicting whether a customer is likely to discontinue telecom services based on historical customer data. Accurate churn prediction enables businesses to take proactive measures to improve customer retention and reduce revenue loss.

\section{Data Description}

The project uses the Telco Customer Churn Dataset obtained from Kaggle. The dataset contains demographic details, service usage information, and payment-related attributes for telecom customers.

\begin{itemize}
    \item \textbf{Total Customers:} 7,043
    \item \textbf{Total Features:} 21
\end{itemize}

\subsection*{Class Distribution}
\begin{itemize}
    \item Customers who did not churn: approximately 73\%
    \item Customers who churned: approximately 27\%
\end{itemize}

\subsection*{Data Preprocessing}
\begin{itemize}
    \item The \texttt{customerID} column was removed as it does not contribute to prediction.
    \item Categorical variables were converted into numerical form using label encoding.
    \item The dataset was split into 80\% training data and 20\% testing data.
\end{itemize}

\section{Exploratory Data Analysis (EDA)}

Exploratory Data Analysis was conducted to understand the patterns and relationships within the dataset. The following insights were observed:

\begin{itemize}
    \item The dataset is imbalanced, with a higher proportion of non-churn customers.
    \item Customers on month-to-month contracts showed higher churn rates compared to long-term contracts.
    \item Customers with shorter tenure were more likely to churn.
\end{itemize}

These insights informed feature understanding and supported the selection of the Random Forest algorithm.

\section{Methodology}

A Random Forest Classifier was chosen for this project due to its ability to handle non-linear relationships and reduce overfitting by combining multiple decision trees.

\begin{itemize}
    \item \textbf{Algorithm Used:} Random Forest Classifier
    \item \textbf{Learning Type:} Supervised Learning (Binary Classification)
    \item \textbf{Library Used:} Scikit-Learn
    \item \textbf{Target Variable:} \texttt{Churn}
\end{itemize}

\subsection*{Model Hyperparameters}

\begin{itemize}
    \item \texttt{n\_estimators = 200}
    \item \texttt{max\_depth = 11}
    \item \texttt{min\_samples\_split = 5}
    \item \texttt{min\_samples\_leaf = 2}
\end{itemize}

\section{Evaluation}

The model was evaluated using classification metrics including accuracy, precision, recall, and F1-score.

\subsection*{Classification Report}

\begin{table}[h]
\centering
\begin{tabular}{lcccc}
\toprule
\textbf{Class} & \textbf{Precision} & \textbf{Recall} & \textbf{F1-Score} & \textbf{Support} \\
\midrule
No Churn (0) & 0.84 & 0.91 & 0.88 & 1036 \\
Churn (1) & 0.68 & 0.52 & 0.59 & 373 \\
\bottomrule
\end{tabular}
\end{table}

\subsection*{Overall Performance Metrics}

\begin{itemize}
    \item \textbf{Accuracy:} 81\%
    \item \textbf{Macro Average F1-Score:} 0.73
    \item \textbf{Weighted Average F1-Score:} 0.80
\end{itemize}

\section{Optimization}

Model performance was improved through hyperparameter tuning of the Random Forest classifier. Adjustments to the number of trees, tree depth, and minimum samples per node helped balance bias and variance, resulting in improved generalization on unseen data.

\section{Team Contribution}

This project was completed as a group effort by \textbf{Kaustubh Hiwanj}, \textbf{Rudraksh Rathod}, and \textbf{Bhargav Patil}. The responsibilities were distributed as follows:

\begin{itemize}
    \item \textbf{Kaustubh Hiwanj:} Model training, evaluation, and preparation of the project presentation (PPT).
    \item \textbf{Rudraksh Rathod:} Development of the project website and creation of the README file.
    \item \textbf{Bhargav Patil:} Documentation and preparation of the project report.
\end{itemize}

All team members actively collaborated throughout the project, properly committed their respective work to the GitHub repository, and contributed collectively to planning, coordination, and successful project completion.

\section{Conclusion}

The Random Forest model achieved over \textbf{80\% accuracy}, demonstrating strong predictive capability for customer churn. The model satisfies the objectives of Capstone Milestone 1 and establishes a solid foundation for future enhancements such as advanced feature engineering and alternative classification algorithms.

\end{document}